\chapter{Status of 21~cm interferometric experiments}

\begin{bf}
  \author{Cathryn M. Trott (ICRAR-Curtin), Jonathan Pober (Brown University)}\\
  
Abstract\\
\end{bf}

This chapter discusses the history, progress, challenges and forecasts for detection and exploration of the spatial structure of the 21~cm brightness temperature signal in the Epoch of Reionisation using interferometric experiments. We discuss GMRT, PAPER, LOFAR, MWA, LEDA, and the future HERA and SKA.


\section{High-level topics}

\begin{itemize}
\item Historical perspective - CMB studies, early arrays
\item Current and recent experiments - description of instruments and experimental design; fundamental differences in analysis and observation methodologies and choices made. LOFAR, MWA, PAPER, GMRT, LEDA (MITEoR?)
\item Exploration of published results; retractions; upper limits; foreground techniques
\item Current challenges: signal loss, leverage, calibration coupling, incomplete sky models and instrument models. How do we know we have detected something?
\item Future prospects for current instruments
\item Future instruments; design, experiments. HERA, SKA. Potential for success. Cross-correlation studies to help with robust signal detection
\end{itemize}



\begin{figure}[]
\begin{center}
\includegraphics[width=0.5\textwidth]{Trott_Pober/01x01-eps-converted-to}
\end{center}
\caption{This is figure 1 in chapter 1.}
\end{figure}

\paragraph{Cras adipiscing} sagittis nunc vel luctus. Suspendisse volutpat augue quis erat semper consequat dignissim tellus euismod. Morbi hendrerit, tellus id aliquam iaculis, nibh leo tincidunt eros, vitae varius ligula felis in mi.

\begin{table}
\caption{Greek Letters.}
\begin{center}
\begin{tabular}{llllllll}
\hline
$\alpha $  & $ \beta $  & $ \gamma $  & $ \delta $  & $ \epsilon $  & $ \varepsilon $  & $ \zeta $  & $ \eta $ \\
 $ \theta $  &  $ \vartheta $  &  $ \gamma $  &  $ \kappa $  &  $ \lambda $  &  $ \mu $  &  $ \nu $  &  $ \xi $ \\
 $ o $  &  $ \pi $  &  $ \varpi $  &  $ \rho $  &  $ \varrho $  &  $ \sigma $  &  $ \varsigma $  &  $$ \\
 $ \tau $  &  $ \upsilon $  &  $ \phi$ &  $ \varphi $  &  $ \chi $  &  $ \psi $  &  $ \omega$  &  $ $ \\
 &  &  &  &  &  &  & \\
$ \Gamma $  & $ \Delta $  & $ \Theta $  &  $ \Lambda $  &  $ \Xi $  &  $ \Pi $  &  $ \Sigma $  & $ \Upsilon $ \\
 $ \Phi$ &  $ \Psi $  &  $ \Omega $  &  &  &  &  &\\
\hline
\end{tabular}
\end{center}\end{table}

\begin{figure}[]
\begin{center}
\includegraphics[width=0.6\textwidth]{Trott_Pober/01x02}
\end{center}
\caption{This is figure 2 in chapter 1.}
\end{figure}


\bibliographystyle{plain}
\bibliography{Trott_Pober/References}


